\section{Experimental Methodology}
\label{sec:experimental-methodology}

This section consolidates the methodological choices used across every
result reported in Sections~\ref{sec:characterization}--\ref{sec:continuation}
so that scope and caveats are stated once rather than scattered.

\textbf{Workloads, capacities, horizons.} All results use the five
families defined in Section~\ref{sec:benchmark-design}
(alibaba-block, metacdn, metakv, twemcache, wiki2018). Offline
target-discriminativeness results (Section~\ref{sec:characterization}) span
all four evaluated capacities (32, 64, 128, 256) and all three horizons
($H\in\{4,8,16\}$). Closed-loop, linkage, mechanistic, and continuation
results (Sections~\ref{sec:closed-loop}--\ref{sec:continuation}) are
reported at capacities 32 and 128 only -- this is the scope of the Tier-1
production closed-loop evaluation on which all three later sections
depend, not a data-availability limitation, and it is stated here once
rather than repeated in each section.

\textbf{Split terminology and two family-specific caveats.} ``test''
denotes a temporally held-out scoring window; ``validation'' denotes a
window used during earlier development that is not temporally held out in
the same sense. Two caveats apply and recur throughout the paper: (i)
\textbf{metacdn's closed-loop evidence is validation-window evidence, not
test-window}, for both of its evaluated capacities; every result that
includes metacdn should be read with this in mind, and Section~\ref{sec:linkage}
reports what changes if metacdn is excluded. (ii) \textbf{metakv is the
only family whose scored window begins at an empty cache} (a cold start,
identifiable from the identity $\mathit{evictions} = \mathit{misses} -
\mathit{capacity}$ holding exactly for metakv and no other family);
Section~\ref{sec:mechanisms} discusses this as a plausible, not
established, contributor to metakv's one policy-ranking exception.

\textbf{Policies and seeds.} Closed-loop replay is evaluated for LRU, MRU,
SIEVE (all deterministic, one run each), and random (20 CRN seeds, seeds
0--19; mean and per-seed standard deviation reported). SIEVE has no
offline counterpart in the labels described in Section~\ref{sec:label-generation}
and is therefore closed-loop-only evidence throughout -- it is not part of
any offline-vs-closed-loop correspondence claim, but is still reported as
a closed-loop reference point (Section~\ref{sec:closed-loop}).

\textbf{Offline metrics.} Per decision: all-tied fraction, unique-winner
fraction, mean/median optimal-set fraction, decision-weighted random-optimal
probability, and random regret (mean and tail percentiles), all computed
per (family, capacity, horizon) stratum before any pooling
(Section~\ref{sec:characterization}).

\textbf{Closed-loop metrics.} Miss ratio per (family, capacity, policy),
and, for random, its across-seed mean and standard deviation; relative
miss-ratio difference vs.\ LRU is reported as the primary comparison
statistic (Section~\ref{sec:closed-loop}).

\textbf{Offline$\leftrightarrow$closed-loop linkage metrics.} Sign
concordance between the offline regret gap and the closed-loop miss-ratio
gap (MRU-vs-LRU and random-vs-LRU only, since these are the two pairs with
both an offline and a closed-loop side), plus Pearson/Spearman/Kendall
correlation between the two gaps' magnitudes, computed per horizon on the
$n=10$ (family, capacity) cells (Section~\ref{sec:linkage}). $n=10$ is
small and is treated as such throughout: correlations are reported with
their exact $n$, never as population-scale claims.

\textbf{Continuation-sensitivity metrics.} Optimal-set Jaccard overlap,
cross-continuation regret (in misses), and strict-reversal fraction,
comparing the canonical LRU-continuation label against MRU- and
mean-random-continuation (10 CRN seeds) relabelings of the same decisions,
on a pre-registered, stratified 5,000-primary + 500-diagnostic decision
sample at capacities 32 and 128, plus a full-population MRU census over
all five families, capacities 32/64/128/256, and horizons 4/8/16
(Section~\ref{sec:continuation}). The primary statistical unit here is
the \emph{decision} (not the candidate row or the continuation seed):
every reported sampled-study statistic is a decision-level cluster
statistic, and the two reported bootstrap confidence intervals (one per
continuation policy; Section~\ref{sec:continuation}) resample at the
decision level; the MRU census is reported as a complete population
enumeration, not an inferential sample.

\textbf{Tie-aware analysis.} Every metric above that could be silently
distorted by ties (optimal-set membership, pairwise preference,
concordance) is computed with ties as a first-class outcome rather than
broken arbitrarily; Section~\ref{sec:background}'s worked example shows
the base case this generalizes from.

\textbf{Provenance and validation strategy.} Every number in this paper
traces to a validated evidence record. The closed-loop, linkage,
mechanistic, and continuation analyses were checked by reproducible scripts
and, where applicable, by independent recomputation of headline values
(Section~\ref{sec:release-validation}). No number in this paper is read
from a partial or unvalidated intermediate run; the full-population MRU
continuation census is used only after complete-record validation and an
independent recheck.
